% Supporting marginal-feasibility lemma, not a novelty or Lean-acceptance claim.
\subsection{Which zero-diagonal graph kernels can be lifted?}

The reduction from an expanded graph to a two-coordinate marginal does not
allow an arbitrary graph kernel to be substituted on the same carrier.
The following finite criterion makes one obstruction explicit.  It is a
weighted fractional clique-decomposition condition, not a new claim about
the general theory of fractional decompositions.

\begin{proposition}[Zero-diagonal marginal criterion]
Let $r\ge2$, let $\Omega$ be finite, and let $\mu(x)>0$ for every
$x\in\Omega$, with $\sum_x\mu(x)=1$.  Let $M:\Omega^2\to[0,\infty)$ be
symmetric with $M(x,x)=0$.  There is a symmetric nonnegative
$U:\Omega^r\to[0,\infty)$ such that
\[
 M(x,y)=\sum_{z_3,\ldots,z_r}
 U(x,y,z_3,\ldots,z_r)\prod_{i=3}^r\mu(z_i)
\]
if and only if there are weights $w_S\ge0$, indexed by the $r$-element
subsets $S\subseteq\Omega$, satisfying
\[
 \mu(x)\mu(y)M(x,y)=\sum_{S\supseteq\{x,y\}}w_S\qquad(x\ne y).
\]
Moreover, one can require $U\le1$ if and only if these equations admit
weights with
\[
 0\le w_S\le(r-2)!\prod_{z\in S}\mu(z).
\]
For either corresponding representation,
$\mathbb E U=r(r-1)\sum_Sw_S$.
\end{proposition}
\begin{proof}
Each diagonal marginal is a finite sum of nonnegative terms with strictly
positive coefficients.  Its vanishing forces $U(x,x,z_3,\ldots,z_r)=0$.
Symmetry therefore makes $U$ vanish on every tuple with a repeated
coordinate.  On an injective tuple with underlying set $S$, its value is
a single number $u_S$, independent of the order of that tuple.  The
remaining $r-2$ coordinates have $(r-2)!$ possible orders, so
\[
 \mu(x)\mu(y)M(x,y)
 =(r-2)!\sum_{S\supseteq\{x,y\}}u_S\prod_{z\in S}\mu(z).
\]
Set $w_S=(r-2)!u_S\prod_{z\in S}\mu(z)$.  Conversely, any weights
satisfying the equations define $U$ to be zero on repeated tuples and
$w_S/((r-2)!\prod_{z\in S}\mu(z))$ on injective tuples.  This kernel is
symmetric, has exactly the stated marginal, and is bounded by one exactly
when the displayed upper bounds hold.  Finally, each $S$ occurs in $r!$
ordered tuples, giving the mean formula.
\end{proof}

The case $r=2$ recovers $U=M$.  If $|\Omega|<r$, the only possible
zero-diagonal marginal is $M=0$.  With zero-mass states, the same statement
applies after restricting to the positive-mass support; no pointwise
conclusion is implied outside that support.

Even the condition that every supported edge belongs to an $r$-clique is
not sufficient.  For $r=3$, put the uniform measure on
$\Omega=\{a,b,c,d\}$, and let $M$ be the adjacency kernel of
$K_4-\{cd\}$.  The only triangles are $abc$ and $abd$.  The $ac$ and $ad$
equations force $w_{abc}=w_{abd}=1/16$, whereas the $ab$ equation requires
$w_{abc}+w_{abd}=1/16$.  Thus no symmetric nonnegative $3$-kernel on this
carrier has marginal $M$, even without an upper bound on $U$.

This criterion is a small constructive bridge rather than a main
contribution.  It turns this fixed-carrier lifting problem into a finite
linear feasibility problem and explains why a graph-level density
counterexample cannot simply be reused as a higher-uniformity marginal.
The classical fractional-decomposition framework is represented, for
example, by Dross, \emph{Fractional Triangle Decompositions in Graphs with
Large Minimum Degree}, SIAM J. Discrete Math. 30 (2016), 36--42,
\url{https://doi.org/10.1137/15M1014310}.  No new dense-graph threshold or
priority claim is made here.
